\begin{abstract}
Numerical solvers used to train PDE surrogates are often treated as neutral label generators, although discretization, governing equations, and output timing can alter apparent model performance. We formulate and evaluate a reference-aware framework for tsunami-like wave emulation in which identical bathymetry--source scenarios are advanced by Hydrostatic, MUSCL-HR, and a study-defined linear weakly dispersive Boussinesq-type reference. Each solver retains its natural time steps while its outputs are interpolated to 50 shared requested times from $8.4$ to $420$ in elapsed benchmark-time units. The source and buffered-grid construction are checked for consistency, while broader numerical verification and independent-comparator studies are reported separately according to their scope. On 2,500 held-out scenarios, global solver-gap RMSEs are $0.00171$, $0.00305$, and $0.00374$ for the Hydrostatic--MUSCL-HR, MUSCL-HR--Boussinesq, and Hydrostatic--Boussinesq pairs, respectively. Among nine Hydrostatic direct models, ConvLSTM has the lowest global relative $L_2$ ($0.203$), followed by U-FNO ($0.224$) and F-FNO ($0.230$). All six cross-reference discrepancy ratios exceed one ($1.10$--$1.72$), so no tested emulator bridges the raw solver gap in this metric. A seven-member FNO ensemble is seed-stable but under-dispersed: test $90\%$ marginal coverage rises from $0.682$ to $0.910$ after validation-only variance scaling, while the high-strength subgroup reaches $0.796$. The synthetic, finite-horizon setting does not establish operational tsunami forecasting or event-scale physical validation.
\end{abstract}
